\documentclass[a4paper,11pt]{ltjsarticle}
\usepackage{amsmath,amssymb}
\usepackage{enumitem}
\usepackage{geometry}
\geometry{margin=25mm}

\title{国政の党勢を地方選挙で補間する状態空間モデル}
\author{}
\date{2026年10月}

\begin{document}
\maketitle

本稿では、党勢の水準と趨勢を、国政選挙と地方選挙の両方から推定する状態空間モデルを提案する。測りたいのは国政での党勢であり、地方選挙は年1回程度しかない国政選挙の間を埋める代理の観測と位置づける。まず先行案(局所回帰、\texttt{local\_regression\_proposal.tex})がこの位置づけを表せなかった理由を述べる(\S\ref{sec:motivation})。次に記号とモデルを定義し(\S\ref{sec:notation}--\ref{sec:model})、パラメータの決め方(\S\ref{sec:params})と出力する量(\S\ref{sec:output})を示す。最後に、先行案との関係(\S\ref{sec:relation})と、検証の方法・未決の点(\S\ref{sec:checks})を述べる。

\section{先行案の問題}
\label{sec:motivation}

先行案の回帰では、国政選挙の扱いを重みでしか調整できず、「国政が本命、地方はその間を埋める」という位置づけを表せなかった。$C_p(T)$と同じ重みを使うと、国政選挙が重みの約9割を占め、傾きはほぼ衆院選2回の差で決まった。この場合、地方選挙は国政選挙の間を埋める役割を果たしていない。一方、誤差の分散を残差から推定して重みにする方法(FGLS)では、国政選挙の重みは2〜5\%に落ちた。2024年と2026年の衆院選の間の大きな変動を、2次式で表せない雑音とみなしたためである。この場合は国政選挙が実質的に無視される。

問題の根は、回帰の式が「党勢そのもの」と「それを観測する選挙」を区別していないことにある。国政選挙の結果は党勢をほぼそのまま表す観測であり、その変動は雑音ではなく信号である。これを表すには、観測されない党勢を状態として持ち、各選挙をその状態の観測として扱うモデルが要る。

\section{記号}
\label{sec:notation}

本節では、モデルを書くのに必要な記号を定義する。選挙イベント$e$、その代表人口$P_e$、投票日$t_e$、無所属を除いて正規化した占有率$\tilde\phi_p(e)$、階層$\ell(e)$は先行案と同じ意味で使う。市区町村長のうち推薦・支持の名簿で未仕分けの選挙を別の階層とすることも先行案と同じである。時刻の単位は月とする。新しく導入する記号は次のとおりである。

\begin{itemize}
  \item $x_p(t)$: 時刻$t$における政党$p$の国政での党勢の水準。衆院選での議席占有率(無所属を除いて正規化)の尺度で測る。直接は観測されない。
  \item $b_p(t)$: 党勢の傾き(月あたりの$x_p$の変化)。直接は観測されない。
  \item $\alpha_{p,\ell}$: 階層$\ell$の選挙で観測される占有率と、国政での党勢$x_p$との水準差。参院選も1つの階層とし、衆院選を基準($\alpha_{p,\text{衆院選}}=0$)とする。
  \item $\lambda_p$: 地方選挙の占有率が、国政での党勢$x_p$の変化にどれだけの比で反応するかを表す負荷。政党ごとに1つで、地方の全階層に共通とする。
  \item $\boldsymbol{\theta}_p(t)$: 状態ベクトル。$\boldsymbol{\theta}_p(t)=(x_p(t), b_p(t), \alpha_{p,\ell_1},\dots,\alpha_{p,\ell_m})^\top$。$\ell_1,\dots,\ell_m$は衆院選以外の階層である。
  \item $\Delta$: 連続する2つの観測時刻の間隔(月)。
\end{itemize}

\section{モデル}
\label{sec:model}

モデルは、党勢がどう動くかを表す状態方程式と、選挙がそれをどう観測するかを表す観測方程式からなる。以下では政党の添字$p$を省く。

\subsection{状態方程式}

党勢の水準と傾きは、どちらも少しずつ不規則に変化するとする。時刻$t$から$t+\Delta$への推移を次の式で表す。
\begin{equation}
  \begin{pmatrix} x(t+\Delta) \\ b(t+\Delta) \end{pmatrix}
  = \begin{pmatrix} 1 & \Delta \\ 0 & 1 \end{pmatrix}
    \begin{pmatrix} x(t) \\ b(t) \end{pmatrix} + \boldsymbol{w},
  \qquad
  \boldsymbol{w}\sim N\!\left(\boldsymbol{0},\;
  q_x\begin{pmatrix}\Delta & 0\\ 0 & 0\end{pmatrix}
  + q_b\begin{pmatrix}\Delta^3/3 & \Delta^2/2\\ \Delta^2/2 & \Delta\end{pmatrix}\right)。
  \label{eq:state}
\end{equation}
$q_x$は水準が不規則に跳ぶ大きさ、$q_b$は傾きが変化する大きさを表す。共分散行列の形は、水準と傾きがそれぞれ連続時間のランダムウォークに従うとしたときに、間隔$\Delta$の推移から導かれるものである(導出は付録\ref{app:Q})。このため観測の間隔が不揃いでもそのまま使える。階層の水準差$\alpha_\ell$は時間によらない定数とする。

\subsection{観測方程式}

選挙イベント$e$は、その投票日の状態を次の式で観測する。
\begin{equation}
  \tilde\phi(e) = \lambda_e\, x(t_e) + \alpha_{\ell(e)} + \varepsilon_e,
  \qquad
  \varepsilon_e \sim N\!\left(0,\, r_e\right)。
  \label{eq:obs}
\end{equation}
負荷$\lambda_e$と観測の誤差の分散$r_e$は、国政選挙と地方選挙で考え方を分ける。

\begin{itemize}
  \item 国政選挙($\ell(e)$が衆院選・参院選): $\lambda_e=1$、$r_e=\sigma_{\text{国}}^2$。$x$は衆院選の議席占有率の尺度で測るので、衆院選の負荷は定義により1である。参院選も同じ国政の議席なので1とし、衆院選との違いは$\alpha_{\text{参院選}}$で表す。国政選挙の結果は、その時点の国政での党勢をほぼそのまま表すとみなし、$\sigma_{\text{国}}$を小さな値に固定する(\S\ref{sec:params})。
  \item 地方選挙: $\lambda_e=\lambda_p$、$r_e=\sigma_\ell^2P_e^{-a}$。小選挙区制のもとでは、衆院の議席占有率は得票の変化を大政党では増幅し、小政党では圧縮する。地方議会の占有率は得票にほぼ比例して動くので、国政の尺度での変化を地方の尺度に直す比$\lambda_p$が要る。地方選挙は自治体ごとの地盤の違いを含むので、誤差は大きい。$\sigma_\ell^2$は階層ごとの大きさ、$a$は人口が大きいほど誤差が小さくなる度合いである。
\end{itemize}

このモデルでは、国政選挙と地方選挙の役割が式の形から決まる。国政選挙があると、誤差が小さいので$x$の推定はその値にほぼ固定される。国政選挙の間は、地方選挙が$\lambda_px+\alpha_\ell$を通じて$x$と$b$を少しずつ更新する。$\alpha_\ell$は、国政選挙と地方選挙が同じ時期に観測された経験から推定される。

\section{パラメータの決め方}
\label{sec:params}

モデルのパラメータは、スタンスとして与えるものとデータから推定するものに分ける。

\begin{itemize}
  \item $\sigma_{\text{国}}$: スタンスとして与える。国政選挙の結果を党勢の定義とみなすので、$\sigma_{\text{国}}=0.005$(占有率で0.5ポイント)とする。0にしないのは数値計算を安定させるためである。
  \item $\lambda_p$: 衆院選の議席と比例代表の得票の関係から、推定の外で決める(\S\ref{sec:seatvote})。国政選挙は政党ごとに4〜6回しかなく、尤度から推定すると識別できない(\S\ref{sec:loading})。
  \item $q_x, q_b, \sigma_\ell^2, a$: 全期間のデータから最尤法で推定する。尤度はカルマンフィルタの予測誤差分解で計算する。
\end{itemize}

国政選挙の誤差をデータから推定しないのは、推定させると先行案のFGLSと同じ問題が起きるからである。国政選挙は2018年以降で衆参合わせて6回しかない。その間の大きな変動を説明するには、$q_x,q_b$を大きくするか$\sigma_{\text{国}}$を大きくするかのどちらかになる。データだけではこの2つを区別しにくく、後者が選ばれれば国政選挙は軽視される。「国政選挙の変動は信号である」というスタンスを、$\sigma_{\text{国}}$を固定することで明示する。

状態の初期値は、政党が初めて占有率を得た投票日(先行案の$E_p(\tau)$と同じ基準)に、平均0、分散の大きい分布として与える。それより前の選挙は使わない。$\alpha_\ell$も同様に、分散の大きい分布から始める。

\section{出力する量}
\label{sec:output}

評価時点$\tau$(各月末)では、$\tau$以前の全ての観測でカルマンフィルタを回し、最後の観測から$\tau$まで状態方程式で予測した値を出力とする。

\begin{itemize}
  \item 党勢変化率$D_p(\tau)=\hat b_p(\tau)$: 月あたりの国政での党勢の変化。標準誤差は状態の共分散行列から得る。
  \item 党勢の現在推定$\hat x_p(\tau)$: 直近の国政選挙の結果を、その後の地方選挙で更新した、国政での党勢の推定値。衆院選の議席占有率(無所属を除く)の尺度で測る。
\end{itemize}

$\hat x_p(\tau)$は$C_p(T)$とは異なる量である。$C_p(T)$は国政・地方の全階層を重み付けして合成した水準であるのに対し、$\hat x_p(\tau)$は「今、衆院選があったら」という国政の尺度での推定である。

\section{先行案との関係}
\label{sec:relation}

先行案の指数重み付き局所多項式回帰は、このモデルの特殊な場合に当たる。観測の間隔が一定で、国政と地方の区別がなく、状態の変化の大きさが一定の比率にある場合、カルマンフィルタの定常状態での推定は、過去の観測に指数的に減衰する重みを付けた当てはめと一致する。先行案の記憶の長さ$h$は、本案では$q_x,q_b$と観測の誤差の比として、データから推定される。

一方、本案で新しく表せるようになったのは次の2点である。

\begin{enumerate}
  \item 国政選挙の結果を、党勢の水準を決める観測として扱える。国政選挙のたびに推定はその結果に合わせ直され、地方選挙はその間を補間する。
  \item 選挙の間隔が不揃いでも、間隔に応じて不確かさが増える。国政選挙から時間が経つほど$\hat x_p(\tau)$の標準誤差は大きくなり、地方選挙の情報の比重が増す。
\end{enumerate}

\section{実データでの検証}
\label{sec:results}

本節では、2018年1月から2026年9月までのデータでモデルを推定し、地方選挙が国政の変化をどの程度先取りできるかを確かめる。検証の対象は、2022年以降の国政選挙(参院選2回、衆院選2回)である。それ以前は水準差$\alpha_\ell$を推定するための国政選挙がまだ少ない。

\subsection{観測に使う階層の条件}

ある政党がほとんど議席を得ない階層は、観測に使わない。具体的には、その階層の選挙のうち占有率が正だった割合が1割未満の場合である。例えばいのちの党の市区町村長選挙では、占有率がほぼ常に0になる。これを正規分布の観測として扱うと、誤差の分散が0と推定され、$\lambda_px+\alpha_\ell$を0に固定する正確な観測になる。その結果、国政選挙が示す党勢の変化を打ち消してしまう。この条件により、いのちの党は地方の階層を全く使わず、国政選挙だけで推定される。

\subsection{国政選挙の予測}
\label{sec:forecast}

地方選挙の情報の価値は、各国政選挙の前日までのデータだけでパラメータを推定し直し、その選挙の結果を予測させて測る。比較の基準は、同じ院の前回の結果をそのまま予測値とする素朴な予測である。名義の変更の影響を受けた2件(2026年衆院選の立憲民主党と公明党)は集計から除き、残る26件で評価した。本小節では地方選挙の負荷を$\lambda_p=1$とし、負荷の決め方は次の小節以降で扱う。

\begin{table}[h]
\centering
\begin{tabular}{lcc}
\hline
予測の方法 & 二乗平均誤差 & 予測の$\pm2$標準偏差に入った割合 \\
\hline
素朴な予測(前回の結果) & 0.083 & --- \\
本モデル($\lambda_p=1$) & 0.064 & 0.65 \\
本モデル+$\alpha_\ell$のランダムウォーク & 0.079 & 0.58 \\
\hline
\end{tabular}
\end{table}

本モデルは素朴な予測より誤差が約23\%小さい。地方選挙は国政の大きな変化の向きを捉えており、大きさの一部を先取りしている。例えば2025年参院選の自民党の後退(19ポイント)のうち6ポイント、2024年衆院選の後退(15ポイント)のうち3ポイントを、選挙前に予測していた。一方、予測の$\pm2$標準偏差に入った割合は0.65で、本来の0.95を大きく下回る。予測の不確かさを過小評価している。

\subsection{地方選挙の反応の大きさ}
\label{sec:loading}

本小節と次の小節では、式\eqref{eq:obs}の負荷$\lambda_e$を地方選挙でどう決めるかを検討する。負荷を1とすると、地方選挙は国政の変化の一部しか先取りしない(\S\ref{sec:forecast})。地方の占有率が、国政の議席占有率ほど大きく動かないためと考えられる。まず負荷を階層ごとに$\lambda_\ell$として尤度から推定すると、推定値は上下限に張り付き、予測は破綻した。各政党の国政選挙は2018年以降で4〜6回しかなく、階層ごとの負荷を識別できない。

次に、全政党・全階層に共通の負荷$\lambda$を固定して比べると、次の結果になった。

\begin{table}[h]
\centering
\begin{tabular}{ccccc}
\hline
$\lambda$ & 対数尤度の合計 & 二乗平均誤差 & $\pm2$標準偏差に入った割合 & 大きな変化の先取り率 \\
\hline
1.00 & 13770 & 0.064 & 0.65 & 0.28 \\
0.70 & 13829 & 0.070 & 0.69 & 0.35 \\
0.50 & 13866 & 0.085 & 0.69 & 0.43 \\
0.35 & 13890 & 0.112 & 0.81 & 0.56 \\
0.20 & 13907 & 0.194 & 0.81 & 0.75 \\
\hline
\end{tabular}
\end{table}

「大きな変化の先取り率」は、前回からの変化が5ポイントを超えた6件について、変化のうち予測で先取りできた割合の平均である。$\lambda$を小さくすると、大きな変化をより多く先取りでき、尤度も高くなる。しかし二乗平均誤差は悪化する。悪化の大半は日本共産党による。共産党は地方議会で強い一方、国政の議席は小選挙区制で圧縮されている。地方での緩やかな退潮を$1/\lambda$倍に増幅すると、国政での大きな後退という誤った予測になる。

この結果は、地方と国政の変化の比が政党の規模に依存することを示している。小選挙区制のもとでは、議席占有率は得票の差を大政党では増幅し、小政党では圧縮する。したがって適切な$\lambda$は、大政党では1より小さく、小政党では1より大きいはずである。共通の$\lambda$では、この両方を同時に満たせない。

\subsection{議席と得票の関係による負荷}
\label{sec:seatvote}

そこで負荷を政党ごとに$\lambda_p$とし、衆院選の議席と得票の関係から、状態空間モデルの推定の外で決めた。地方議会の占有率は得票にほぼ比例して動き、衆院の議席は得票の変化を増幅・圧縮して動く、という見立てである。衆院選$j$での政党$p$の議席占有率(無所属を除いて正規化)を$x_{pj}$、比例代表の得票率を$v_{pj}$とし、3乗比の法則の一般形
\begin{equation}
  \operatorname{logit} x_{pj} = k\,\operatorname{logit} v_{pj} + c
\end{equation}
を全政党・全選挙をまとめて最小二乗法で当てはめる。$k$は議席の得票に対する弾力性である。弾力性$k$は全政党で共通とし、切片は政党ごとに異なってよいとすると、政党$p$の直近の衆院選での傾きは$dx/dv=k\,x(1-x)/\{v(1-v)\}$となる。負荷はその逆数
\begin{equation}
  \lambda_p = \frac{v_{pj}(1-v_{pj})}{k\,x_{pj}(1-x_{pj})}
\end{equation}
とした。ここで$j$は推定時点より前の直近の衆院選である。比例代表の得票率を総務省の集計表から読み取れなかった政党は、共通の曲線上にあるとみなして$v$を$x$から逆算する。

2017年以降の4回の衆院選で当てはめると$k=1.89$となった。2026年9月末時点の$\lambda_p$は、自由民主党0.57、立憲民主党0.40、日本維新の会0.58、国民民主党0.81、公明党1.02、参政党1.15、日本共産党2.58、いのちの党6.90である。大政党で1より小さく、小政党で1より大きいという、\S\ref{sec:results}の結果から予想した形になっている。

この$\lambda_p$で国政選挙の予測を検証し直すと、二乗平均誤差は0.057となり、素朴な予測(0.083)より約31\%小さくなった。検証では、各国政選挙の前日より前の衆院選だけから$k$と$\lambda_p$を決め直している。$\lambda=1$との比較は次のとおりである。
\begin{itemize}
  \item 対数尤度の合計は13770から13857に上がった。
  \item 日本共産党の二乗平均誤差は0.072から0.031に下がり、共通の$\lambda$で生じた問題は解消した。
  \item 参政党は0.053から0.071に悪化した。議席3の2024年衆院選から$\lambda_p=2.53$を決めたため、2025年参院選の躍進を地方選挙から先取りする分が小さくなった。
  \item $\pm2$標準偏差に入った割合は0.65で変わらず、大きな変化の先取り率は0.28から0.31にわずかに上がった。
\end{itemize}
議席が数議席しかない政党では$x_{pj}$が整数の議席数の揺らぎに大きく左右され、$\lambda_p$が不安定になる。いのちの党の6.90(2026年衆院選の1議席による)はその例である。

\subsection{傾きの動き方}
\label{sec:slope}

本小節では、式\eqref{eq:state}のように傾きをランダムウォークさせる必要があるかを確かめる。全期間で推定した$q_b$は、全政党で$10^{-8}$以下となった。$q_b=10^{-9}$でも、傾きが12か月で変わる大きさは$\sqrt{12q_b}\approx10^{-4}$ポイント/月にすぎない。尤度は、傾きがほとんど変化しないモデルを選んでいる。

そこで、傾きの動き方が異なる次の4つのモデルを比べた。負荷は\S\ref{sec:seatvote}の$\lambda_p$とする。
\begin{itemize}
  \item 傾きなし: $b\equiv0$。$x$はランダムウォークだけで動く。
  \item 傾き一定: $q_b=0$。$b$は全期間で1つの値をとり、データから推定される。
  \item ランダムウォーク: 式\eqref{eq:state}のとおり。
  \item 減衰: $db=-\beta b\,dt+dW_b$。傾きは一時的な勢いを表し、減衰率$\beta$(月あたり)で0に戻る。$\beta$も最尤法で推定する。
\end{itemize}

\begin{table}[h]
\centering
\begin{tabular}{lcccc}
\hline
傾きの動き方 & 対数尤度の合計 & 二乗平均誤差 & $\pm2$標準偏差に入った割合 & 大きな変化の先取り率 \\
\hline
傾きなし & 13889 & 0.059 & 0.54 & 0.27 \\
傾き一定 & 13857 & 0.057 & 0.62 & 0.30 \\
ランダムウォーク & 13857 & 0.057 & 0.65 & 0.31 \\
減衰 & 13889 & 0.059 & 0.54 & 0.27 \\
\hline
\end{tabular}
\end{table}

4つのモデルは、実質的に「傾きなし」と「傾き一定」の2つに分かれた。ランダムウォークは$q_b$がほぼ0なので傾き一定と一致する。減衰では、推定した$\beta$が多くの政党で1〜5(上限)となり、勢いは1か月程度で消えるので、傾きなしと一致する。

傾き一定でも、推定された傾きはどの政党でも0と区別できない。2026年9月末の傾きと標準誤差は、例えば自由民主党で$+0.15\pm0.62$、参政党で$+0.22\pm0.15$、立憲民主党で$+0.03\pm0.41$ポイント/月である。対数尤度の合計は傾きなしの方が31高いが、これは傾きの事前分布の扱いによる差である。散漫な事前分布(分散0.01)から傾きを推定すると、事前分散と事後分散の比の対数の半分が尤度から差し引かれる。政党ごとの対数尤度の差はこの量とほぼ一致する(日本維新の会で6.5対6.4、自由民主党で2.8対2.8)。データへの当てはまりは、2つのモデルで変わらない。国政選挙の予測では、傾き一定の方が二乗平均誤差がわずかに小さく、被覆も良いが、26件での差であり決め手にはならない。

以上から、2018年以降のデータでは、党勢の傾きは検出できない。党勢は、なめらかな趨勢としてよりも、水準の不規則な跳び($q_x$)として動いている。先行案の検証で、傾きのない局所回帰が3か月先の予測で最良だったことも、この結果と整合する。

\subsection{水準の跳び}
\label{sec:jump}

本小節では、水準の変化を正規分布のランダムウォークではなく跳びとして表すモデルを試す。きっかけは、国政選挙の予測誤差の形である。傾きなしのモデルで、予測誤差を予測の標準偏差で割った値は、標準偏差3.6、超過尖度2.75となった。モデルが正しければそれぞれ1と0である。雑音の大きさ$q_x$は数千件の地方選挙で決まり、地方選挙は党勢がふだんあまり動かないことを示す。一方で国政選挙はときどき大きな変化を示す。正規分布のランダムウォークは、この2つを1つの分散で表そうとして両立できない。

そこで式\eqref{eq:state}の$x$に、月あたりの発生率$\rho$で分散$J$の跳びを加えた。間隔$\Delta$のうちに1回以上跳ぶ確率は$1-e^{-\rho\Delta}$で、2回以上跳ぶ場合も1回の跳びで近似する。フィルタは「跳ばなかった」「跳んだ」の2成分の正規混合で状態を持ち、次の時刻に進む前に、平均と分散が等しい1つの正規分布にまとめ直す。$\rho$と$J$は他のパラメータと同時に最尤法で推定する。

\begin{table}[h]
\centering
\begin{tabular}{lcccc}
\hline
モデル & 対数尤度の合計 & 二乗平均誤差 & $\pm2$標準偏差に入った割合 & 標準化誤差の標準偏差 \\
\hline
傾きなし & 13889 & 0.059 & 0.54 & 3.6 \\
傾きなし+跳び & 13941 & 0.058 & 0.65 & 3.3 \\
傾き一定 & 13857 & 0.057 & 0.62 & 3.4 \\
傾き一定+跳び & 13912 & 0.056 & 0.73 & 3.1 \\
\hline
\end{tabular}
\end{table}

跳びを加えると、対数尤度は50余り上がり、被覆も改善した。追加したパラメータは政党あたり2つ(8政党で16)なので、尤度の改善はこれを上回る。しかし、改善は部分的にとどまる。
\begin{itemize}
  \item 尤度の改善の大半は日本維新の会(35)と立憲民主党(13)によるもので、他の政党ではほぼ0である。日本維新の会の跳びの標準偏差は上限の30ポイントに張り付いており、推定は安定していない。
  \item 大きな変化を示した6件では、予測値も予測の標準偏差も跳びのないモデルとほとんど変わらない。それぞれの選挙の前日までのデータには、その政党にこの大きさの跳びがまだ現れていないので、跳びの分散が小さく推定される。
  \item 標準化誤差の標準偏差は3.1〜3.3で、1からはまだ遠い。
\end{itemize}

最大の予測誤差を調べると、外れ値は国政選挙だけでなく地方選挙の側にもある。最大級の外れ値は、無所属が議席の大半を占める町村議会の選挙で生じる。無所属を除いて正規化すると、わずかな政党議席が大きな占有率になるからである。例えば2021年の三朝町議会では、立憲民主党の議席は8\%だが、無所属(92\%)を除くと占有率は1になる。日本維新の会でも、2023年の有田市議会(議席13\%、無所属73\%)が0.48になる。こうした選挙の予測誤差は、予測の標準偏差の8〜16倍に達する。地方選挙の誤差は正規分布から大きく外れており、そのために$\sigma_\ell^2$が大きく推定され、他の地方選挙の情報も割り引かれている可能性がある。状態の跳びだけでは、この問題は扱えない。

\subsection{地方選挙の観測誤差}
\label{sec:localerr}

本小節では、地方選挙の観測誤差の形を変える2つの方法を試す。1つ目は、確率$\pi$で誤差の分散が$\kappa$倍になる外れ値の成分を混ぜる方法である。2つ目は、無所属を除いた議席数$n^\ast_e$での二項分布の分散を誤差に加える方法である。後者では、地方選挙の誤差の分散を
\begin{equation}
  r_e = \sigma_\ell^2P_e^{-a} + \frac{\hat p_e(1-\hat p_e)}{n^\ast_e}
\end{equation}
とする。$\hat p_e$はその選挙の占有率の予測値である。三朝町議会の例では$n^\ast_e=1$なので、占有率が1でも情報の少ない観測として扱われる。

\begin{table}[h]
\centering
\begin{tabular}{lcccc}
\hline
モデル & 二乗平均誤差 & $\pm2$標準偏差に入った割合 & 標準化誤差の標準偏差 & 大きな変化の先取り率 \\
\hline
傾きなし & 0.059 & 0.54 & 3.6 & 0.27 \\
傾きなし+外れ値の成分 & 0.077 & 0.46 & 11.0 & 0.07 \\
傾きなし+二項分布の分散 & 0.058 & 0.65 & 5.0 & 0.31 \\
傾き一定+二項分布の分散 & 0.055 & 0.65 & 3.2 & 0.38 \\
傾きなし+二項分布の分散+跳び & 0.055 & 0.62 & 3.5 & 0.37 \\
\hline
\end{tabular}
\end{table}

外れ値の成分を混ぜると、推定は退化した。参政党の市区町村議会選挙では、85\%の選挙で占有率が0である。最尤推定は、正規の成分の標準偏差を0.25ポイントまで小さくしてこの0の山に当て、正の占有率をすべて外れ値の成分に回した($\pi=0.15$、$\kappa$は上限の$10^4$)。その結果、$\lambda_px+\alpha_\ell$は0に固定され、$q_x$はほぼ0になって党勢が動かなくなった。対数尤度は大きく上がるが、0の山に幅の狭い正規分布を当てたことによる見かけの改善であり、国政選挙の予測は大きく悪化した。

二項分布の分散を加えると、国政選挙の予測はわずかに改善した。傾き一定、または跳びと組み合わせた場合の二乗平均誤差0.055と先取り率0.37〜0.38は、これまでで最も良い。跳びを加えたときの対数尤度の上昇は、二項分布の分散がある場合の方が大きい(52に対して275)。ただし標準化誤差の標準偏差は3〜5で、不確かさの過小評価は残る。対数尤度は、二項分布の分散が下限として効くので2000ほど下がるが、0を多く含むデータに正規分布を当てた尤度どうしの比較には意味が乏しい。

2つの結果は、同じ原因を指している。小政党の地方の占有率は0を多く含み、正規分布ではそもそも表せない。観測誤差の分布を手直ししても、この点は変わらない。根本的には、議席数を二項分布などで直接モデル化し、状態を$\operatorname{logit}x$の空間に置く必要がある。\S\ref{sec:seatvote}の議席と得票の関係もlogitの空間では直線になるので、負荷の扱いも簡単になる。

\section{logitの状態と議席数の観測}
\label{sec:glm}

本節では、\S\ref{sec:localerr}の結論に従い、状態を$\operatorname{logit}x$の空間に置き、地方選挙を議席数の分布で観測するモデルを作る。あわせて、自治体ごとの地盤の差と、参院選の扱いを見直す。実装は\texttt{state\_space\_glm.py}である。

\subsection{モデル}

状態は$z=\operatorname{logit}x$と階層の水準差$\alpha_\ell$とする。\S\ref{sec:slope}の結果から傾きは置かない。$z$は\S\ref{sec:jump}までと同じくランダムウォーク(分散$q_x$)に従う。

地方選挙$e$では、無所属を除いた議席数を$n^\ast_e$、そのうち政党の議席数を$k_e$とし、
\begin{equation}
  k_e \sim \text{BetaBin}(n^\ast_e,\ \mu_e,\ \phi_\ell),\qquad
  \operatorname{logit}\mu_e = \eta_e = \gamma_p\,z(t_e) + \alpha_{\ell(e)} + u_{m(e)}
  \label{eq:glmobs}
\end{equation}
とする。ベータ二項分布は、選挙ごとの真の占有率が平均$\mu_e$、精度$\phi_\ell$のベータ分布でばらつき、議席がそこから二項分布で決まるという2段のモデルである。首長選($n^\ast_e=1$)ではベルヌーイ分布になる。0や1の占有率も自然に扱え、無所属が大半の議会のように$n^\ast_e$が小さい選挙は、情報の少ない観測になる。推薦が複数政党に割れた首長選では$k_e$が0.5などの割合になるので、ベルヌーイ分布の尤度を割合のまま使う。

負荷$\gamma_p$は、正規分布のモデルの$\lambda_p$(\S\ref{sec:seatvote})をlogitの空間に直して決める。$\lambda_p$は「地方の占有率$\mu$の変化は得票率$v$の変化に等しい」という見立てだった。$d\mu=dv$と衆院選の議席と得票の関係$\operatorname{logit}x=k_{\text{衆}}\operatorname{logit}v+c_{\text{衆}}$から
\begin{equation}
  \gamma_p = \frac{d\operatorname{logit}\mu}{dz} = \frac{v_p(1-v_p)}{k_{\text{衆}}\,\bar\mu_p(1-\bar\mu_p)}
  \label{eq:gammap}
\end{equation}
となる。$v_p$は衆参いずれかの直近の国政選挙での比例代表の得票率、$\bar\mu_p$は推定時点より前の地方議会の議席の合計から求めた占有率である。どちらもデータから計算でき、推定の外で決まる。地方の占有率が得票率より低い政党ほど$\gamma_p$は大きく、地方での変化は国政に小さく写される。$\bar\mu_p=v_p$なら$\gamma_p=1/k_{\text{衆}}\approx0.53$で、全政党で共通になる。$x$は必ず0と1の間に入る。

国政選挙は、\S\ref{sec:params}のスタンスのまま$z$をほぼ正確に観測する。議席数に0.5を足してlogitをとり、誤差$\sigma_{\text{国}}$をlogitの尺度に換算する。観測が状態の非線形関数なので、地方選挙ごとに$\eta_e$についての1次元のラプラス近似で事後分布を正規分布に直し、同じ事後分布を与える正規の擬似観測で状態を更新する。対数尤度も同じ近似で求める。

\subsection{自治体ごとの水準}

式\eqref{eq:glmobs}の$u_m$は自治体$m$に固有の水準で、平均0、分散$\tau_\ell^2$の確率変数とする。同じ自治体の前回の選挙と今回の占有率の相関は、市区町村議会で自由民主党0.83、日本共産党0.83、公明党0.81と高い。各党の占有率の自治体間のばらつきの多くは、選挙をまたいで残る地盤の差である。$u_m$は状態ベクトルには入れず、自治体ごとに事後分布の平均と分散を覚えておき、次にその自治体の選挙があったときの事前分布にする($u_m$と状態との相関は無視する)。推定するパラメータは$\tau_\ell^2$だけで、階層ごとに1つである。

自治体の選挙は4年に1度だが、データは2018年から始まるので、2022年以降の地方選挙にはすべて前回の選挙がある。実質的には、同じ自治体の前回からの変化を観測に使うことになる。仮想的な議会で全国を集計する$C_p$とは違い、各自治体は別々に動き、その共通の動きを$z$として取り出す。

\subsection{参院選の扱い}

\S\ref{sec:model}では、参院選を$\alpha_{\text{参院選}}$だけずれた$x$の観測とした。しかし参院選は各党1〜2回しかなく、$\alpha_{\text{参院選}}$は決まらない。そこで、参院選の改選議席にも議席と得票の関係$\operatorname{logit}x_{\text{参}}=k_{\text{参}}\operatorname{logit}v+c_{\text{参}}$を当てはめ、得票率$v$を介して衆院とつなぐ。2016年以降の4回の参院選では$k_{\text{参}}=1.59$、残差の標準偏差は0.26であった。$v$を消去すると
\begin{equation}
  \operatorname{logit}x_{\text{参}} = \frac{k_{\text{参}}}{k_{\text{衆}}}\,z + c_{\text{参}} - \frac{k_{\text{参}}}{k_{\text{衆}}}c_{\text{衆}}
\end{equation}
となる。参院選の観測の負荷を$k_{\text{参}}/k_{\text{衆}}$とし、$\alpha_{\text{参院選}}$の事前分布の平均を右辺の定数項、分散を2つの曲線の残差分散から決める。どちらも推定の外で決まる。

\subsection{結果}

これまでと同じ国政選挙の予測26件で比べた結果を示す。

\begin{table}[h]
\centering
\begin{tabular}{lcccc}
\hline
モデル & 二乗平均誤差 & 予測区間に入った割合 & 標準化誤差の標準偏差 & 先取り率 \\
\hline
正規分布、傾きなし(\S\ref{sec:slope}) & 0.059 & 0.54 & 3.6 & 0.27 \\
本節、自治体の効果なし & 0.102 & 0.77 & 1.7 & 0.44 \\
本節、自治体の効果あり & 0.086 & 0.77 & 2.0 & 0.43 \\
本節、自治体の効果+参院選のつなぎ & 0.083 & 0.81 & 2.2 & 0.40 \\
\hline
\end{tabular}
\end{table}

不確かさの見積もりは大きく改善した。予測区間($\pm2$標準偏差)に入る割合は0.81、標準化誤差の標準偏差は2.2で、正規分布のモデルよりかなり現実に近い。自治体の効果を入れると対数尤度は1080上がり、ベータ二項分布の精度$\phi_\ell$はほぼ二項分布の値まで上がった。地方の過分散の大半は、自治体の地盤の差だったことになる。参院選のつなぎによる改善は小さいが、決まらない$\alpha_{\text{参院選}}$を推定の外の情報で補えるので採用する。

一方、二乗平均誤差は正規分布のモデルより大きい。誤差の中央値は素朴な予測と同程度で、大外れが少数ある。主なものは次の2件である。
\begin{itemize}
  \item 日本維新の会の2024年衆院選: 実際0.085に対し予測0.205。2023年の統一地方選挙で地方の議席が大きく増え、推定した党勢は0.3前後まで上がったが、国政の議席は増えなかった。
  \item 国民民主党の2026年衆院選: 実際0.061に対し予測0.265。2025年以降の地方での伸びが、国政の大きな伸びと解釈された。
\end{itemize}
この表の結果は、負荷を全政党に共通の$\gamma$とし、尤度から決めたものである(\S\ref{sec:constants})。どちらも、小政党が地方で議席を大きく伸ばした場面である。logitの空間では、地方の占有率の比での変化が、国政では$1/\gamma$倍に増幅される。$\gamma=0.53$なら、地方で議席が倍になった政党は、国政で4倍近い議席を得ると予測される。小選挙区では議席が得票より速く伸びるので、この増幅自体は不自然ではない。2件の外れは、それまでの地方と国政の関係から外れることが起きたと読める。モデルの出力は「地方選挙から見た国政の党勢」の推定であり、外れが起きれば次の国政選挙で党勢はその結果に合わせ直される。ただし、同じ外れが地方で伸びる小政党で毎回起きるなら、それは系統的な偏りである。今後の国政選挙で、予測の外れを政党の規模や地方での伸びと並べて確かめる。正規分布のモデルは占有率の差で地方と国政をつないでいたので、この増幅が起きなかった。

\subsection{定数の扱い}
\label{sec:constants}

本小節では、モデルに含まれる定数を、データから推定するもの、推定の外の別データから決めるもの、結果に効かないことを確かめたものに分ける。マジックナンバーに依存するモデルは採らない、という方針による。

\begin{itemize}
  \item 負荷: 尤度ではほとんど決まらない。全政党に共通の$\gamma$として推定すると、全期間のデータで、$\gamma$を0.3、0.4、0.53、0.7、0.9に固定して8政党の対数尤度の合計を比べると、0.3〜0.53の差は2.4以内でほぼ平らであり、0.7で4.7、0.9で15.2低かった(公明党は最適化が局所解に落ちたので除いた。以後、最適化は複数の初期値から始める)。$\gamma$の95\%の範囲はおよそ0.3〜0.6である。式\eqref{eq:gammap}の$\gamma_p$を使っても、8政党の対数尤度の合計は共通の$\gamma=0.45$と0.7しか違わない。尤度で区別できないので、推定の外のデータから計算できる$\gamma_p$を採る。
  \item 衆参の議席と得票の曲線と、参院選の負荷・水準差の事前分布: 総務省の公表値から計算する。
  \item 国政選挙の誤差$\sigma_{\text{国}}$: 0.0025、0.005、0.01と変えると、2026年9月末の$\hat x$の変化は1ポイント以下であった。月次の系列でも、国民民主党(最大4.6ポイント)を除けば2.3ポイント以内に収まる。尤度はこの範囲でほぼ平らで、データからは決まらない。0.02まで大きくすると、国民民主党の系列が最大12ポイント動く。結果に効かない範囲の中央の0.005を使う。
  \item 状態の事前分散: 10を100にしても、月次の系列の変化は2.2ポイント以内であった。この程度の差は、最適化の収束の違いでも生じる。
  \item 階層を使う条件(占有率が正の選挙が1割以上、10件以上): ベータ二項分布は0の多いデータも扱えるので、条件を外す。外しても、主要政党の系列の変化は自由民主党で2.3ポイント、立憲民主党で0.6ポイントであった。大きく変わるのは国政選挙の議席がごく少ない政党で、いのちの党の最新値は0.020から0.001になり、日本共産党は最初の国政選挙より前の2019年前半の値が大きく変わった。議席が数議席の政党では、国政選挙の誤差$\sigma_{\text{国}}$をlogitに換算すると大きくなり、国政選挙が党勢をつなぎ止められないためである。
  \item 国政選挙のlogitで足す0.5議席: 議席0の政党でlogitを有限にするための補正で、議席の少ない政党にしか効かない。
\end{itemize}

以上の扱いで、国政選挙の予測を検証し直した。比較のため、共通の$\gamma$を各選挙の前日までのデータで0.3、0.45、0.6の中から尤度で選び直した場合も示す。選ばれた$\gamma$は、2022年参院選から2025年参院選までが0.6、2026年衆院選が0.45であった。$\gamma_p$も各選挙の前日までのデータで計算する。2026年9月末時点の$\gamma_p$は、日本共産党0.13、公明党0.21、自由民主党0.55、立憲民主党0.76、日本維新の会1.32、参政党1.66、国民民主党1.69、いのちの党2.27である。

\begin{table}[h]
\centering
\begin{tabular}{lcccc}
\hline
政党 & 素朴な予測 & 正規分布(\S\ref{sec:slope}) & 本節、共通の$\gamma$ & 本節、$\gamma_p$ \\
\hline
自由民主党 & 0.158 & 0.102 & 0.110 & 0.108 \\
立憲民主党 & 0.060 & 0.042 & 0.040 & 0.037 \\
日本共産党 & 0.013 & 0.028 & 0.007 & 0.008 \\
国民民主党 & 0.040 & 0.034 & 0.084 & 0.038 \\
日本維新の会 & 0.016 & 0.016 & 0.053 & 0.022 \\
公明党 & 0.024 & 0.027 & 0.034 & 0.038 \\
参政党 & 0.071 & 0.057 & 0.075 & 0.074 \\
\hline
\end{tabular}
\caption{国政選挙の予測の絶対誤差の平均(政党別)。いのちの党は省略した。}
\end{table}

共通の$\gamma$では、国民民主党、日本維新の会、公明党で素朴な予測より大きく外れた。地方での伸び縮みを国政に写す際の増幅が、中規模の政党で過大になっている。外れが2件の出来事にとどまらず中規模の政党に共通して現れるので、\S\ref{sec:glm}の結果の項で述べた「系統的な偏り」に当たる。$\gamma_p$にすると、国民民主党は0.084から0.038、日本維新の会は0.053から0.022になり、この偏りはほぼ解消した。他の政党はほとんど変わらない。

\begin{table}[h]
\centering
\begin{tabular}{lccc}
\hline
 & 二乗平均誤差 & 誤差の中央値 & 予測区間に入った割合 \\
\hline
素朴な予測 & 0.083 & 0.021 & --- \\
正規分布(\S\ref{sec:slope}) & 0.059 & 0.024 & 0.54 \\
本節、共通の$\gamma$ & 0.093 & 0.029 & 0.81 \\
本節、$\gamma_p$ & 0.073 & 0.028 & 0.73 \\
\hline
\end{tabular}
\caption{国政選挙の予測26件の誤差。}
\end{table}

26件全体では、$\gamma_p$で二乗平均誤差は素朴な予測を下回ったが、誤差の中央値は及ばず、正規分布のモデルにも及ばない。予測区間に入る割合は0.73に下がったが、正規分布のモデルの0.54よりは良い。残る大きな外れは、自由民主党の2026年衆院選(実際0.686、予測0.385)と参政党の2025年参院選(実際0.124、予測0.014)である。どちらも直前の地方選挙に兆しのない国政の大きな変化で、地方選挙を代理にする限り予測できない。

\section{説明のための推定}
\label{sec:explain}

本節では、モデルの目的を国政選挙の予測から党勢の動きの説明に置き直し、出力と検証をそれに合わせる。国政選挙の結果を事前に当てることは、そもそも難しい。予測と実際の結果の差は、その選挙が党勢に与えた衝撃そのものである。説明として問うべきことは、国政選挙と国政選挙の間の党勢の動きを地方選挙でどれだけ説明できるか、推定された動きがどの選挙によるものか読み取れるか、の2点である。

\subsection{平滑化}

過去の時点の党勢は、その後の選挙も使って推定する。前向きのフィルタの結果から、固定区間の平滑化(RTS)で各時点の状態の平均と分散を求める。地方選挙はフィルタでのラプラス近似の擬似観測のまま扱い、自治体の効果$u_m$は平滑化しない。月次の系列は平滑化の値で出力する。最新の月はフィルタの値と一致する。

\subsection{国政選挙の補間}

国政選挙を1回ずつ除き、残りの全期間のデータ(前後両側)から平滑化でその結果を再現できるかを確かめる。除いた選挙の議席と得票の点は、$\gamma_p$と参院選のつなぎの計算からも除く。比較の基準は、同じモデルから地方選挙を除いた、国政選挙だけの補間である。地方選挙を使うことで何が上乗せされるかが、これで直接分かる。各党の最初の国政選挙を除く34件で比べた。

\begin{table}[h]
\centering
\begin{tabular}{lcccc}
\hline
 & 二乗平均誤差 & 誤差の中央値 & 予測区間に入った割合 & 標準化誤差の標準偏差 \\
\hline
国政選挙だけ & 0.078 & 0.023 & 0.74 & 1.6 \\
正規分布、傾きなし(\S\ref{sec:slope}) & 0.069 & 0.023 & 0.74 & 2.9 \\
地方選挙も使う & 0.069 & 0.021 & 0.79 & 2.1 \\
地方選挙も使う、$z_N$あり(\S\ref{sec:natdev}) & 0.080 & 0.023 & 0.82 & 1.7 \\
\hline
\end{tabular}
\caption{国政選挙の補間の誤差(34件)。}
\end{table}

地方選挙を使うと、誤差はすべての指標で小さくなる。改善が大きいのは、日本維新の会(絶対誤差の平均0.023から0.016)、参政党(0.064から0.050)、自由民主党(0.159から0.140)である。地方選挙は、国政選挙の間の党勢の動きの一部を説明している。ただし改善の幅は大きくない。残る大きな外れは、自由民主党の2026年衆院選(実際0.686、補間0.397)、参政党の2025年参院選(0.124、0.026)、国民民主党の2025年参院選(0.150、0.073)で、いずれも地方選挙に現れなかった国政の動きである。

\subsection{選挙ごとの寄与}

ある期間の党勢の推定値の変化は、その期間の選挙ごとの寄与にちょうど分けられる。傾きを置かないので、前向きのフィルタでは予測の段階で$z$の平均は動かず、$z$は観測ごとの更新でだけ動くからである。選挙$e$の寄与は、その選挙での期待される占有率$\operatorname{logit}^{-1}(\gamma_p m_z+\alpha_\ell+u_m)$と実際の占有率の差に、利得を掛けたものになる。月次の出力には、その月の選挙ごとに、期待される占有率、実際の占有率、$z$の変化を載せる。予測と実際の差が衝撃であるという見方を、選挙単位で表したものである。

\subsection{水準の跳び}

\S\ref{sec:jump}の跳びをこのモデルに入れても、全期間の対数尤度の上がりは各党1.6以下で、増える2つのパラメータに見合わない。跳びは採らない。

\subsection{国政選挙にだけ現れる部分}
\label{sec:natdev}

補間の標準化誤差の標準偏差は2.1で、国政選挙の結果は地方選挙から期待されるより大きく動く。そこで国政の党勢を$z+z_N$とし、地方選挙の負荷は$z$にだけかける。$z_N$は始点で0、分散$q_N$のランダムウォークで、$q_N$は尤度から推定する。党勢の変化を、地方選挙にも現れる部分と国政選挙にだけ現れる部分に分けられる。

全期間の対数尤度は、自由民主党で7.1、国民民主党で9.5、参政党で7.7、立憲民主党で2.9上がり、他の党ではほぼ変わらない(パラメータは各党1つ増える)。推定された$q_N$(年あたり、logitの尺度)は、参政党0.53、自由民主党0.058、国民民主党0.027で、地方選挙の$q_x$より1桁以上大きい。これらの党の国政での動きの多くは、地方選挙に現れない。一方、補間では予測区間に入る割合が0.82、標準化誤差の標準偏差が1.7と不確かさの見積もりは良くなったが、二乗平均誤差は0.080に悪化した。悪化の大半は参政党で、$q_N$が大きいため、除いた選挙の値を地方選挙ではなく前後の国政選挙から決めるようになる。$z_N$は本体には入れず、診断用の出力とする。

\subsection{相乗りの首長選}

相乗りの首長選を推薦した政党で等分する代わりに、選挙直前の各党の当選確率$\hat\mu_{pe}$に比例して分ける。1段目は等分のデータで推定し、その$\hat\mu_{pe}$で分け直したデータで推定し直す(\texttt{trend\_model.tex}の相乗りの節)。補間(34件)の二乗平均誤差は0.067から0.068と変わらない。分け直しは等分の不自然さを除くために採る。

その後、首長選そのものを観測から外したので(次節)、この分け直しは使っていない。

\subsection{首長選を使わない形}

地方選挙を議会選挙(都道府県議会、市区町村議会)だけにする形も、同じ補間(36件)で比べた。二乗平均誤差は、すべての階層を使う形で0.081、市区町村長選を除くと0.078、知事選も除いて議会だけにすると0.078、市区町村議会だけにすると0.080である。首長選を除いても補間は悪くならず、知事選はほとんど効いていない。首長選の党派は推薦・支持に頼るしかなく、市区町村長選の推薦・支持を載せる全国首長名簿は選挙の9か月〜1年9か月後にしか出ないので、首長選は使わない。都道府県議会は補間では区別できないが、名簿の遅れの問題がないので使う。詳細は\texttt{trend\_model.tex}の地方選挙の範囲の節にある。

\subsection{全政党をまとめたモデル}

全政党の議席を1つの多項分布で観測するモデルも試した。状態は政党ごとの$z_p$と$\alpha_{p\ell}$を並べたもので、首長選は「当選者は推薦した政党のどれかの陣営である」という観測(組の尤度)として扱える。パラメータは政党ごとのモデルの推定値を使い、$\phi_{p\ell}$は無限大とした。補間(34件)の二乗平均誤差は0.070で、政党ごとのモデル(0.068)より少し悪く、$\pm2$標準偏差に入る割合も0.74にとどまる。$\phi_{p\ell}$を落としたぶん地方選挙を信じすぎ、区間が狭くなるためと考えられる。

まとめたモデルでは、政党どうしの$z_p$の増分に相関$R$を入れられる。$R$を、政党ごとのモデルの平滑化した$z_p$の月ごとの増分の相関行列$R_{\text{emp}}$から計算し、$(1-\lambda)I+\lambda R_{\text{emp}}$の$\lambda$を尤度で決めると、$\lambda=1$で全期間の対数尤度が10.3上がる。全ペアに共通の相関$\rho$では、尤度はほとんど上がらない(最大0.5)。補間では、$\lambda$と$R_{\text{emp}}$を除いた後のデータで決め直しても、二乗平均誤差は0.070から0.069に下がるだけで、政党ごとのモデルには届かない。$R_{\text{emp}}$は、期間の短い政党(いのちの党)を含む組で不安定でもある。まとめたモデルは採らない。

\section{検証と未決の点}
\label{sec:checks}

モデルが妥当かどうかは、標準化した予測誤差(イノベーション)で確かめる。モデルが正しければ、各観測の予測誤差をその予測分散の平方根で割った値は、平均0・分散1で、時間方向に相関しない。階層ごと、および統一地方選挙とそれ以外に分けて確認する。特に重要なのは、国政選挙の直前の予測誤差である。地方選挙から補間した$\hat x_p$が、次の国政選挙の結果をどの程度予測できていたかが、地方選挙を代理に使うことの妥当性を直接示す。ただし国政選挙は6回しかないので、この確認は粗いものにとどまる。

次の点は未決であり、検証の結果を見て判断する。

\begin{itemize}
  \item 出力する量。傾きが検出できない(\S\ref{sec:slope})ので、党勢変化率$D_p=\hat b_p$は、その時点の勢いではなく全期間の平均の傾きに近い値になり、ほぼ常に0と区別できない。変化を出すなら、$\hat x_p(\tau)$と$\hat x_p(\tau-12)$の差のように、水準の推定値の差として定義し直す方がよい。
  \item 議席の少ない政党の$\lambda_p$の扱い(\S\ref{sec:seatvote})。上限を設ける、複数回の衆院選で平均する、などが考えられる。
  \item 予測の不確かさの過小評価(\S\ref{sec:jump})。状態の跳びを加えても、予測の$\pm2$標準偏差に入る割合は0.65〜0.73にとどまる。地方選挙の誤差の手直し(\S\ref{sec:localerr})でも解消しなかったが、logitの状態と議席数の観測(\S\ref{sec:glm})で予測区間に入る割合は0.81まで改善した。
  \item 水準差$\alpha_\ell$を定数とするか、ゆっくり変化する状態とするか。ランダムウォークにした場合は予測が悪化した(\S\ref{sec:results})。
  \item 名義の変更。立憲民主党と中道改革連合の名義の問題は、国政の観測としてそのまま入る。「段差は党勢低下と読む」という既存の方針に従い、立憲民主党の党勢の大きな低下として現れる。
\end{itemize}

\appendix
\section{状態雑音の共分散行列の導出}
\label{app:Q}

式\eqref{eq:state}の共分散行列は、水準と傾きのそれぞれに連続時間の雑音が入るモデルを、間隔$\Delta$だけ解いたものである。ここでは、左上の成分が$\Delta^3/6$ではなく$\Delta^3/3$になる理由を含めて導出する。

連続時間のモデルを
\begin{equation}
  dx = b\,dt + dW_x, \qquad db = dW_b
\end{equation}
とする。$W_x$と$W_b$は互いに独立なブラウン運動で、単位時間あたりの分散はそれぞれ$q_x$、$q_b$である。傾き$b$はランダムウォークし、水準$x$はその$b$を積分したものに、自身のランダムウォークを加えたものである。

時刻$0$から$\Delta$まで解くと、次のようになる。
\begin{align}
  b(\Delta) &= b(0) + \int_0^\Delta dW_b(u), \\
  x(\Delta) &= x(0) + b(0)\,\Delta + \int_0^\Delta W_b(s)\,ds + W_x(\Delta)。
\end{align}
第2式の積分は、積分の順序を入れ替えると
\begin{equation}
  \int_0^\Delta W_b(s)\,ds = \int_0^\Delta (\Delta-u)\,dW_b(u)
\end{equation}
と書ける。時刻$u$に傾きが受けた変化はその後も残るので、時刻$\Delta$までに水準を$(\Delta-u)$だけずらす、という意味である。

独立な増分の和の分散は、各増分の係数の2乗を積分したものになる(伊藤の等長性)。したがって
\begin{align}
  \mathrm{Var}[x(\Delta)] &= q_x\Delta + q_b\int_0^\Delta (\Delta-u)^2\,du = q_x\Delta + q_b\frac{\Delta^3}{3}, \\
  \mathrm{Cov}[x(\Delta), b(\Delta)] &= q_b\int_0^\Delta (\Delta-u)\,du = q_b\frac{\Delta^2}{2}, \\
  \mathrm{Var}[b(\Delta)] &= q_b\int_0^\Delta du = q_b\Delta
\end{align}
となり、式\eqref{eq:state}の行列が得られる。ここでの分散と共分散は、時刻$0$の状態を与えたときの条件付きのものである。

左上の成分の$1/3$は、$(\Delta-u)^2$を1回積分して出る係数である。ある量を3回積分したときの$1/3!=1/6$とは由来が違う。$1/6$は、例えば一定の加速度による変位を時間で積分したときのように、ずれそのもの(平均)に現れる係数である。分散では、ずれ$(\Delta-u)$を先に2乗してから積分するので、$1/3$になる。

離散時間のモデルで「各ステップの間は傾きの変化率が一定」とすると、$\boldsymbol{g}=(\Delta^2/2,\ \Delta)^\top$として共分散行列は$q\,\boldsymbol{g}\boldsymbol{g}^\top$となり、成分は$\Delta^4/4$、$\Delta^3/2$、$\Delta^2$である。この形は、間隔の異なる2つのステップを続けて進めた結果が、まとめて1ステップ進めた結果と一致しない。選挙は不等間隔に行われるので、本稿では連続時間から導いた形を使う。連続時間の形では、間隔$\Delta_1$と$\Delta_2$の推移を続けて進めると、間隔$\Delta_1+\Delta_2$の推移と同じ分布になる。

\end{document}
